% ═══════════════════════════════════════════════════════════════
% LEHRERHINWEISE - DS-01: Mi nueva casa
% Klasse 9, Unidad 2
% ═══════════════════════════════════════════════════════════════

\documentclass[11pt, a4paper]{article}

\usepackage[utf8]{inputenc}
\usepackage[T1]{fontenc}
\usepackage[ngerman]{babel}

\usepackage{helvet}
\renewcommand{\familydefault}{\sfdefault}

\usepackage[margin=2cm, top=2.5cm, bottom=2cm]{geometry}
\usepackage{enumitem}
\usepackage{tabularx}
\usepackage{booktabs}
\usepackage{array}
\usepackage{multicol}

\usepackage{fancyhdr}
\usepackage{lastpage}

\usepackage{xcolor}
\usepackage{tcolorbox}
\tcbuselibrary{skins}

% ═══════════════════════════════════════════════════════════════
% FARBEN
% ═══════════════════════════════════════════════════════════════

\definecolor{akzent}{HTML}{C41E3A}
\definecolor{hellgrau}{HTML}{F5F5F5}
\definecolor{mittelgrau}{HTML}{666666}
\definecolor{basis}{HTML}{4CAF50}
\definecolor{standard}{HTML}{2196F3}
\definecolor{erweiterung}{HTML}{9C27B0}

% ═══════════════════════════════════════════════════════════════
% BOXEN
% ═══════════════════════════════════════════════════════════════

\newtcolorbox{infobox}[1][]{
    colback=hellgrau, colframe=akzent, boxrule=1pt, arc=0pt,
    left=10pt, right=10pt, top=8pt, bottom=8pt,
    title={#1}, fonttitle=\bfseries, coltitle=white, colbacktitle=akzent
}

\newtcolorbox{diffbox}[2][]{
    colback=white, colframe=#2, boxrule=1pt, arc=0pt,
    left=8pt, right=8pt, top=6pt, bottom=6pt,
    title={#1}, fonttitle=\bfseries\small, coltitle=white, colbacktitle=#2
}

% ═══════════════════════════════════════════════════════════════
% KOPF- UND FUSSZEILE
% ═══════════════════════════════════════════════════════════════

\pagestyle{fancy}
\fancyhf{}
\renewcommand{\headrulewidth}{0.5pt}
\renewcommand{\footrulewidth}{0pt}
\lhead{\footnotesize\textsc{Lehrerhinweise} · DS-01}
\rhead{\footnotesize Mi nueva casa · Klasse 9}
\cfoot{\footnotesize\thepage\,/\,\pageref{LastPage}}

\setlength{\parindent}{0pt}
\setlength{\parskip}{6pt}

% ═══════════════════════════════════════════════════════════════
% DOKUMENT
% ═══════════════════════════════════════════════════════════════

\begin{document}

% ═══════════════════════════════════════════════════════════════
% TITEL
% ═══════════════════════════════════════════════════════════════

\begin{center}
    {\LARGE\bfseries Lehrerhinweise: Mi nueva casa}\\[6pt]
    {\large DS-01 · Klasse 9 · Unidad 2}
\end{center}

\vspace{8pt}

% ═══════════════════════════════════════════════════════════════
% ÜBERBLICK
% ═══════════════════════════════════════════════════════════════

\begin{infobox}[Stundenübersicht]
\begin{tabular}{@{}ll@{}}
\textbf{Thema:} & Mi nueva casa -- Nuevos contenidos \\
\textbf{Dauer:} & 90 Minuten (Doppelstunde) \\
\textbf{Materialien:} & AB-01, LB-01, LP-01, SLIDES-01, Beamer \\
\textbf{Quelle:} & Qué pasa 2, Unidad 2 (S. 26-27, 31-33) \\
\end{tabular}
\end{infobox}

\vspace{8pt}

% ═══════════════════════════════════════════════════════════════
% FEINZIELE
% ═══════════════════════════════════════════════════════════════

\section*{Feinziele}

\begin{itemize}[leftmargin=1.5em]
    \item \textbf{FZ1:} SuS \textbf{benennen} 7 Hausräume + 17 Möbel/Gegenstände (AFB I)
    \item \textbf{FZ2:} SuS \textbf{wenden} Bewegungsverben mit Präpositionen an (AFB II)
    \item \textbf{FZ3:} SuS \textbf{ersetzen} direkte Objekte durch Pronomen (AFB II)
    \item \textbf{FZ4:} SuS \textbf{beschreiben} einen Umzug mit Verben + Pronomen (AFB II/III)
\end{itemize}

% ═══════════════════════════════════════════════════════════════
% STUNDENVERLAUF
% ═══════════════════════════════════════════════════════════════

\section*{Stundenverlauf}

\begin{tabularx}{\textwidth}{@{}c|l|X|c|l@{}}
\textbf{Min} & \textbf{Phase} & \textbf{Aktivität} & \textbf{SF} & \textbf{Material} \\
\midrule
10 & Einstieg & Bild Wohnung zeigen, Vorwissen aktivieren (hay, ser, estar) & PL & SLIDES 1-2 \\
\midrule
20 & Input I & Vokabeln: Hausräume + Möbel präsentieren (Audio, Bildkarten) & PL & SLIDES 3-4 \\
& & $\rightarrow$ AB Kasten V1 (LESEN) & & LB-01 \\
\midrule
15 & Üben I & Zuordnungsübung: Bild → Wort & EA & AB Aufg. 1 \\
& & Selbstkontrolle mit LB & & \\
\midrule
15 & Input II & Bewegungsverben + Präpositionen einführen (Symbole) & PL & SLIDES 5-6 \\
& & $\rightarrow$ AB Kasten G1 (AUSFÜLLEN) & & \\
\midrule
10 & Üben II & Lückentext: Verben + Präpositionen & EA & AB Aufg. 2 \\
\midrule
10 & Input III & Objektpronomen einführen (lo, la, los, las) & PL & SLIDES 7-8 \\
& & $\rightarrow$ AB Kasten G2 (AUSFÜLLEN) & & \\
\midrule
10 & Anwendung & Partnerarbeit: Umzug beschreiben (mit Pronomen) & PA & AB Aufg. 3-4 \\
\midrule
& & $\rightarrow$ AB Kasten R1 (LESEN) & & \\
\bottomrule
\end{tabularx}

% ═══════════════════════════════════════════════════════════════
% DIFFERENZIERUNG
% ═══════════════════════════════════════════════════════════════

\section*{Differenzierung}

\begin{multicols}{3}

\begin{diffbox}[{[B] Basis}]{basis}
\footnotesize
\begin{itemize}[leftmargin=1em, itemsep=2pt]
    \item Nur 10 Kernvokabeln
    \item Vokabelliste mit Bildern
    \item Aufgabe 4: nur 2 Sätze
    \item Mehr Zeit bei Übungen
\end{itemize}
\end{diffbox}

\columnbreak

\begin{diffbox}[{[S] Standard}]{standard}
\footnotesize
\begin{itemize}[leftmargin=1em, itemsep=2pt]
    \item Alle 24 Vokabeln
    \item Angeleitete Übungen
    \item Aufgabe 4: vollständig
    \item Normale Zeitvorgaben
\end{itemize}
\end{diffbox}

\columnbreak

\begin{diffbox}[{[E] Erweiterung}]{erweiterung}
\footnotesize
\begin{itemize}[leftmargin=1em, itemsep=2pt]
    \item Zusatz: Eigene Texte
    \item ``Mi habitación ideal''
    \item Komplexere Dialoge
    \item Peer-Tutoring
\end{itemize}
\end{diffbox}

\end{multicols}

% ═══════════════════════════════════════════════════════════════
% LÖSUNGEN
% ═══════════════════════════════════════════════════════════════

\section*{Lösungen}

\textbf{Kasten G1 (Präpositionen):}
\begin{itemize}[leftmargin=1.5em, itemsep=0pt]
    \item bajar \textbf{de / del} · subir \textbf{a / al} · poner \textbf{en}
    \item sacar \textbf{de / del} · traer \textbf{de / del} · llevar \textbf{a / al}
\end{itemize}

\textbf{Kasten G2 (Pronomen):}
\begin{itemize}[leftmargin=1.5em, itemsep=0pt]
    \item Singular: \textbf{lo} (mask.) / \textbf{la} (fem.)
    \item Plural: \textbf{los} (mask.) / \textbf{las} (fem.)
\end{itemize}

\textbf{Aufgabe 1:} 1-b, 2-d, 3-f, 4-c, 5-a, 6-e, 7-j, 8-i, 9-g, 10-h, 11-l, 12-k

\textbf{Aufgabe 2:}
\begin{enumerate}[label=\alph*), leftmargin=1.5em, itemsep=0pt]
    \item bajan ... del
    \item suben ... al
    \item pone ... en
    \item saca ... de
\end{enumerate}

\textbf{Aufgabe 3:}
\begin{enumerate}[label=\alph*), leftmargin=1.5em, itemsep=0pt]
    \item La ponemos en el salón.
    \item Los subimos al piso.
    \item Lo llevamos al dormitorio.
    \item Las sacamos del camión.
    \item La bajamos del piso.
    \item Lo traemos del salón.
\end{enumerate}

% ═══════════════════════════════════════════════════════════════
% TYPISCHE FEHLER
% ═══════════════════════════════════════════════════════════════

\section*{Typische Fehler}

\begin{tabularx}{\textwidth}{@{}lXX@{}}
\textbf{Fehler} & \textbf{Beispiel} & \textbf{Korrektur} \\
\midrule
Falsche Präposition & *La ponemos \textit{a} el salón. & La ponemos \textbf{en} el salón. \\
de + el nicht kontrahiert & *Bajo \textit{de el} camión. & Bajo \textbf{del} camión. \\
Falsches Genus Pronomen & *\textit{Lo} mesa ponemos... & \textbf{La} mesa ponemos... \\
Pronomen nach Verb & *Ponemos \textit{la} en el salón. & \textbf{La} ponemos en el salón. \\
\bottomrule
\end{tabularx}

% ═══════════════════════════════════════════════════════════════
% HINWEISE
% ═══════════════════════════════════════════════════════════════

\section*{Didaktische Hinweise}

\begin{itemize}[leftmargin=1.5em]
    \item \textbf{Einstieg:} Reales Umzugsszenario nutzen -- viele SuS haben Erfahrung damit
    \item \textbf{Vokabeln:} Audio aus QP2 nutzen für Aussprache (CD Track X)
    \item \textbf{Bewegungsverben:} Symbole ($\uparrow\downarrow\oplus\ominus\leftarrow\rightarrow$) helfen bei der Memorierung
    \item \textbf{Objektpronomen:} Betonung auf \textbf{Stellung vor Verb} -- häufigster Fehler!
    \item \textbf{Partnerarbeit:} Rollenwechsel nach 4 Minuten einplanen
\end{itemize}

% ═══════════════════════════════════════════════════════════════
% AUSBLICK
% ═══════════════════════════════════════════════════════════════

\section*{Ausblick}

\textbf{DS-02 (Martes 13):} Wiederholung der Inhalte + Verbindung mit pretérito perfecto aus Unidad 1

\textbf{DS-03 (Lunes 19):} Klausurvorbereitung -- alle Themen U1 + U2

\textbf{Klassenarbeit:} 19./20. Januar

\end{document}
